\chapter*{Carrot Cake}
\addcontentsline{toc}{chapter}{Carrot Cake}
\markboth{Carrot Cake}{Carrot Cake}

% TAGS: gâteaux, carottes, noisettes
\textbf{Tags:} \texttt{gâteaux}, \texttt{carottes}, \texttt{noisettes}

% SERVINGS: 6 à 8
% PREP_TIME: 20 min
% COOK_TIME: 1 hour
% BEGIN_METADATA
\begin{center}
Portions : 6 à 8 \quad Prep time: 20 min \quad Cook time: 1 hour
\end{center}
% END_METADATA

\section*{Ingrédients}

\begin{multicols}{2}
\begin{itemize}[leftmargin=*]
\item \textbf{375 g} \ingredient{carottes}
\item \textbf{185 g} \ingredient{poudre de noisette}
\item \textbf{150 g} \ingredient{beurre}
\item \textbf{3} \ingredient{œufs}
\item \ingredient{oudres}
\item \textbf{1 c. à café} \ingredient{cannelle en poudre}
\item \textbf{75 g} \ingredient{farine}
\item \textbf{1 sachet} \ingredient{levure chimique}
\item \textbf{150 g} \ingredient{cassonade}
\item \textbf{1 pincée} \ingredient{sel fin}
\item \textbf{30 g} \ingredient{beurre pour le moule}
\item \textbf{50 g} \ingredient{cassonade pour le moule}
\item \textbf{100 g} \ingredient{beurre mou}
\item \textbf{250 g} \ingredient{Philadelphia}
\item \textbf{100 g} \ingredient{sucre glace}
\end{itemize}
\end{multicols}

\section*{Ustensiles}

\begin{itemize}[leftmargin=*]
\item \cookware{poêle}
\item \cookware{moule}
\item \cookware{feuille de papier de cuisson}
\item \cookware{couteau}
\item \cookware{fourchette}
\item \cookware{fouet}
\end{itemize}

\section*{Préparation}

\subsection*{Préparation du gâteau}

\begin{enumerate}
\item Pelez, lavez puis râpez les \ingredient{carottes}. Grillez la moitié de la \ingredient{poudre de noisette}, à sec dans une \cookware{poêle}. Laissez-la ensuite refroidir. Faites fondre doucement le \ingredient{beurre}.
\item Préchauffez le four à 210 \textdegree{}C. Fouettez les \ingredient{œufs} en omelette puis mélangez-les avec les carottes, les 2 \ingredient{oudres} de noisette (grillée et non grillée) et la \ingredient{cannelle en poudre}.
\item Ajoutez ensuite le beurre fondu, la \ingredient{farine} et la \ingredient{levure chimique} tamisées, la \ingredient{cassonade} et le \ingredient{sel fin}.
\item Beurrez l'intérieur du \cookware{moule} haut (16cm x 7.5cm) avec du \ingredient{beurre pour le moule} puis poudrez-le avec de la \ingredient{cassonade pour le moule}. Versez la pâte et enfournez pour \timer{1 heure}. Si le gâteau dore trop rapidement, baissez légèrement le thermostat du four et déposez sur le gâteau une \cookware{feuille de papier de cuisson}. Vérifiez la cuisson en enfonçant dans le gâteau la lame d'un \cookware{couteau} qui doit ressortir sèche. Patientez \timer{5 mn} et démoulez-le.
\end{enumerate}

\subsection*{Glaçage}

\begin{enumerate}
\item Préparez le glaçage. Travaillez le \ingredient{beurre mou} à la \cookware{fourchette} ou au \cookware{fouet} puis incorporez peu à peu le \ingredient{Philadelphia}. Ajoutez le \ingredient{sucre glace} sans cesser de travailler le mélange.
\item Réfrigérez le glaçage \timer{15 mn}, puis déposez-le en couche épaisse sur le carrot cake froid.
\end{enumerate}

\chapter*{Gateau au yaourt}
\addcontentsline{toc}{chapter}{Gateau au yaourt}
\markboth{Gateau au yaourt}{Gateau au yaourt}

\section*{Ingrédients}

\begin{multicols}{2}
\begin{itemize}[leftmargin=*]
\item \textbf{5} \ingredient{citrons bio}
\item \textbf{100 g} \ingredient{beurre demi-sel}
\item \textbf{3} \ingredient{œufs}
\item \textbf{150 g} \ingredient{sucre glace}
\item \textbf{200 g} \ingredient{farine}
\item \textbf{0.5 sachet} \ingredient{levure chimique}
\item \textbf{1 pot} \ingredient{yaourt grec}
\item \textbf{2 cs} \ingredient{graines de pavot}
\item \textbf{1 cs} \ingredient{jus de citron}
\item \textbf{100 g} \ingredient{sucre glace}
\item \textbf{120 ml} \ingredient{eau}
\item \textbf{220 g} \ingredient{sucre en poudre}
\end{itemize}
\end{multicols}

\section*{Ustensiles}

\begin{itemize}[leftmargin=*]
\item \cookware{jatte}
\item \cookware{moule}
\item \cookware{papier cuisson}
\item \cookware{grille}
\item \cookware{casserole}
\item \cookware{papier cuisson}
\end{itemize}

\section*{Préparation}

\subsection*{La pate}

\begin{enumerate}
\item Préchauffez le four à 180\textdegree{}C. 
\item Préparez la pâte : lavez et séchez 3 \ingredient{citrons bio}. Râpez finement leur zeste et pressez les fruits. Faites fondre \ingredient{beurre demi-sel} et laissez-le tiédir.
\item Dans une \cookware{jatte}, fouettez \ingredient{œufs} avec le \ingredient{sucre glace}. 
\item Ajoutez la \ingredient{farine}, la \ingredient{levure chimique} et le beurre fondu. Incorporez \ingredient{yaourt grec}. Versez le jus et les zestes de citron, ajoutez \ingredient{graines de pavot} et mélangez la pâte. Versez la préparation dans un \cookware{moule} à cake chemisé de \cookware{papier cuisson}. Enfournez \timer{40 à 45 min} : le cake doit être moelleux à l'intérieur et doré à l'extérieur. 
\item Démoulez le cake et laissez-le refroidir sur une \cookware{grille}.
\end{enumerate}

\subsection*{Glaçage}

\begin{enumerate}
\item Préparez le glaçage en mélangeant le \ingredient{jus de citron} et le \ingredient{sucre glace}. Nappez le cake refroidi de ce glaçage et laissez-le figer à température ambiante.
\end{enumerate}

\subsection*{Les citrons confits}

\begin{enumerate}
\item Lavez et séchez les 2 citrons bio restants. Coupez-les en fines rondelles. Dans une \cookware{casserole}, préparez un sirop avec \ingredient{eau} et \ingredient{sucre en poudre}. Portez à ébullition, incorporez les rondelles de citrons et laissez confire \timer{15 min}. Laissez refroidir sur du \cookware{papier cuisson} puis décorez le cake.
\end{enumerate}

